\section{Introduction}

The data products from NSF-DOE Vera C. Rubin Observatory (funded by the U.S. National Science Foundation and the U.S. Department of Energy's Office of Science) and its Legacy Survey of Space and Time (LSST; \cite{2019ApJ...873..111I}) are expected to be accessed and analyzed by ten thousand scientists and students worldwide.
Supporting this large and diverse community to access the novel LSST data set requires an innovative and proactive community science model.
This paper provides an overview of the Rubin Observatory's community science model, including policies, guidelines, and workflows for documentation, tutorials, activities, user support, and issue resolution.
An evaluation of the model's current implementation and how these components are working "in the wild", with an eye towards the challenges of scaling up while maintaining inclusion and access and deriving metrics for scientific performance, are also included.


\section{Context}\label{sec:context}

At the time this manuscript was prepared in May 2026, Rubin was in its "Early Science" era.
Only a small amount of data -- relative to the future expected datasets -- had been released.

\subsection{Early Science Program and Data Previews}

Early science is defined as any science enabled by Rubin for its community prior to the first annual data release, Data Release 1 \citedsp{RTN-011}.
It includes early alert production, incremental template generation, and the Data Previews based on commissioning and science validation data.
As of May 2026 two Data Previews had been released:
Data Preview 0 (DP0\footnote{\url{https://dp0.lsst.io/}}), based on simulated LSST-like images and catalogs,
and Data Preview 1 \citep[DP1\footnote{\url{https://dp1.lsst.io/}}][]{2026arXiv260323786V}, based on observations with the LSST Commissioning Camera.
Alerts had started streaming to brokers as of February 2026, with the release of Prompt Products and Data Preview 2 (DPS) from the LSST Camera expected later in 2026.

\subsection{The Rubin Science Platform (RSP)}

The Rubin Science Platform (RSP\footnote{\url{https://rsp.lsst.io/}}) provides authenticated access to the LSST data products, tools, and services \citep{LSE-319, LDM-542, 2024ASPC..535..227O}.
It offers a JupyterLab python environment with the Rubin software installed and maintained; a Portal aspect with a graphical user interface for data exploration and visualization; and an API aspect for remote data access.
A limited amount of storage space and compute resources are available to each account, with additional resources available through the Resource Allocation Committee \citep{RTN-084}.
In this manuscript, the term "user" refers both generally to a user of the Rubin data products and also specifically users of the RSP.
The RSP is available to Rubin data-rights holders only \citedsp{RDO-013}.
At the time of writing in May 2026, there were several thousand user accounts initiated in the RSP.

\subsection{The Community Science Team (CST)}

Situated within the Data Management department, the CST is composed of $\sim$10 staff members, each at 0.5 to 1.0 full-time equivalent (FTE) hours.
They are a mix of postdoctoral fellows and junior and senior research staff, with research expertise across the Rubin science pillars: understanding the nature of dark matter and dark energy; taking an inventory of the Solar System; mapping the Milky Way; and exploring variable and transient phenomena.
The CST is responsible for developing and implementing the Community Science Model, and for moderating the Rubin Community Forum and managing user support requests.

\subsection{Use cases and user profiles}

The development of Rubin's Community Science Model began with defining the RSP user profiles and use cases that the model would serve \citep{RTN-002}.
User profiles include students, professional scientists (the occasional, moderate, and heavy user), and users with disabilities (e.g., blind, deaf, and neurodivergent users).
The use cases are specific examples of simple to advanced analyses that any user might need to do with the RSP, involving data query, retrieval, visualization, and analysis.
These user profiles and use cases motivated the development of the Community Science Model.

\subsection{The DP0 Delegates program}

In 2021, at the time DP0 was released, RSP access was limited to 300 user accounts.
Due to this limit imposing an application-and-selection process, the DP0 Delegates program was established \citedsp{RTN-004}.
The selection process prioritized inclusion, with the goal of assembling an early set of users who represented the diversity of the broader science community, who could provide useful feedback to Rubin and also share what they learned back with their communities (i.e., "seeding expertise").
Delegates, in turn, benefited from early access and an accelerated learning experience.
The program ended up being only slightly oversubscribed, and 300 additional user accounts were added yearly until the RSP could support open registration for any data-rights holder.
Early experiences serving the Delegates strongly motivated the development of the Community Science Model.


\section{Community Science Model}\label{sec:model}

Providing user support at scale starts with preventing issues from needing staff attention in the first place, by providing sufficient documentation and learning resources, and by enabling peer-to-peer and crowd-sourced solutions.
When issues do rise to the level of requiring staff intervention, the priority is to provide a positive user experience and a timely and consistent solution by following established internal workflows.

\subsection{Documentation}

The Rubin Observatory For Scientists website\footnote{\url{https://rubinobservatory.org/for-scientists}} is designed as the first point of contact.
The page hierarchy is "shallow", with only three layers: the main landing page, "hub" topical pages, and "content" pages.
This keeps the menu easily navigable and information discoverable.
Content pages are high-level descriptions that link users to technical documentation, and typically end in a link to the Rubin Community Forum.

Technical documentation, which includes data release documentation websites and Rubin Tech Notes, is served from  \url{https://www.lsst.io/}.
All technical documentation adheres to the Rubin user documentation style guide\footnote{\url{https://developer.lsst.io/user-docs/index.html}} which follows the industry practice of topic-based writing ("every page is page one").
This design is also intended to keep detailed information easily discoverable.
Metadata and schema (catalog column names and descriptions\footnote{\url{https://sdm-schemas.lsst.io/}}) are essential technical documentation, particularly important because they are delivered to users with the data itself, and are thus the most likely to be consulted.

The two most challenging aspects of developing user-facing documentation for the Rubin data products are one,
the LSST science pipelines \citep{10.71929/rubin/2570545} are still in active development and are constantly being improved,
and two, the complexity of the image processing and cataloged measurements that the software produces.
These challenges are mitigated by the documentation being open-source and version-controlled in GitHub,
enabling all staff members to contribute.
Keeping documentation accessible and up-to-date is an essential component of a scalable user support model.


\subsection{Tutorials}

Guidelines in RTN-045...



\subsection{Learning opportunities}

Live interactive learning opportunities, both virtual and in-person, have the dual benefits of educating and engaging users.

\begin{enumerate}
\item Rubin Science Assemblies (weekly, virtual): Held Thursday mornings, the weeks alternate between practical, hands-on tutorials and open drop-in "office hour" sessions for live user support.
\item Rubin Data Academy (annual, virtual): Typically the week after the summer meeting of the American Astronomical Society, with multiple sessions a day serving global timezones. A week of interactive tutorials and networking sessions to provide an accelerated learning experience for anyone, but especially summer students starting Rubin projects. 
\item Rubin Community Workshop (annual, hybrid): Held in the middle of northern summer to bring together Rubin users and staff for a week of invited and contributed scientific presentations and discussion.
\item Support for regional/topical meetings and custom seminars: The CST aims to fulfill all user requests for talks and tutorials or find an available Rubin staff member. Resources are limited, with preference given to small and underserved institutes and virtual talk requests. When Rubin staff travel they're encouraged (and supported by the CST) to give talks and tutorials at local institutes.
\end{enumerate}

Virtual events scale easily to hundreds of participants following along with presentations and tutorials, and the ease of recording and hosting videos in, e.g., YouTube, extends the reach even further.


\subsection{User support and issue resolution}

Guidelines for Forum RTN-097


\subsection{Accessibility}

External accessibility reviews...

Colorblind-friendly filter palette... x-rubin effort.

Supporting science at small and underserved institutes...

Translations...

\subsection{Supporting the LSST Science Collaborations}

Helping the helpers, training the trainers.
Liaison program.


\section{Performance}\label{sec:performance}

Metrics remain a challenge.
Got number of papers, number of users, but... 

The SciX library and tech note...

Annual reports...

Forum status ... 

User surveys...

Users Committee...



\section{Future Work}\label{sec:future}

AI, rubin bot, integrations

More short video tutorials, especially for the Portal.

Class modules, so undergraduate professors can incorporate training and undergrads can get credit for learning to analyze Rubin data eg in prep for summer projects (or non summer).